\begin{trapbox}
	\begin{description}[style=nextline, leftmargin=2.8em, labelsep=0.5em, itemsep=0.35em]
		\item[\textbf{\textcolor{CTrap}{易错点一（连乘范围）：}}] 混淆累乘的上下限，常将 $\prod_{k=1}^{n-1}$ 误写为 $\prod_{k=1}^{n}$ 或 $\prod_{k=0}^{n-1}$，导致多乘或少乘因子。
		\item[\textbf{\textcolor{CTrap}{正确理解（正确范围）：}}] 公式中连乘从 $k=1$ 到 $n-1$，共 $n-1$ 个因子；对应递推式依次用了 $n=1,2,\dots,n-1$。
		\item[\textbf{\textcolor{CTrap}{错因分析（迭代混淆）：}}] 误认为递推式写了 $n$ 次，实际只写出了 $n-1$ 个等式。
	\end{description}

	\medskip
	\begin{description}[style=nextline, leftmargin=2.8em, labelsep=0.5em, itemsep=0.35em]
		\item[\textbf{\textcolor{CTrap}{易错点二（初始验证）：}}] 直接使用通项公式计算 $a_1$，但公式要求 $n\ge 2$，导致 $n=1$ 时公式无定义或结果可疑。
		\item[\textbf{\textcolor{CTrap}{正确理解（先验首项）：}}] 公式仅在 $n\ge 2$ 时成立；对于 $n=1$ 应单独给出 $a_1$ 并验证与公式是否一致。
		\item[\textbf{\textcolor{CTrap}{错因分析（忽略范围）：}}] 不注意定理中的条件 $n\ge 2$，机械代入 $n=1$ 产生错误。
	\end{description}

	\medskip
	\begin{description}[style=nextline, leftmargin=2.8em, labelsep=0.5em, itemsep=0.35em]
		\item[\textbf{\textcolor{CTrap}{易错点三（非零条件）：}}] 在变形 $\frac{a_{n+1}}{a_n}=f(n)$ 时，默认 $a_n\neq 0$ 但未预先验证，一旦出现零项会导致错误。
		\item[\textbf{\textcolor{CTrap}{正确理解（确保非零）：}}] 使用累乘法前必须确认数列所有项均不为零；若题目未明确，需加检验。
		\item[\textbf{\textcolor{CTrap}{错因分析（分母疏忽）：}}] 忽略除法对分母非零的要求，错误地认为递推式总能做除法。
	\end{description}
\end{trapbox}
